\section{A real auxiliary function with balanced coefficients}\label{sec:balanced}

The wave construction will use a bounded real trigonometric polynomial
whose coefficient magnitudes are close to the square roots of its
frequency weights.  These weights must have the form $|a\cdot v|$, where
$a$ is a Fourier frequency and $v$ is one common vector.  We adapt the auxiliary-torus construction of
\cite[Section~2]{OpenAILowerEnvelope2026}.  The additional step is to
round the coefficients of each spreading factor to a common magnitude.

\begin{proposition}\label{prop:balanced}
For every $0<\delta<1/100$, there are an integer $m\ge2$, a real
trigonometric polynomial $F$ on $\T^m$, and $v\in\R^m$ with the
following properties.  Its Fourier support
$\mathcal A\subset\Z^m\setminus\{0\}$ consists of at least two pairs
$\{a,-a\}$, with distinct pairs nonparallel.  If
\[
 c_a=\widehat F(a),\qquad \lambda_a=a\cdot v,\qquad w_a=|\lambda_a|,
\]
then, for an absolute constant $C$,
\begin{equation}\label{eq:balanced}
 \begin{gathered}
 \norm F_\infty\le1+\delta,\qquad w_a>0,\\
 1-C\delta\le\sum_{a\in\mathcal A}w_a<1,\qquad
 1\le\frac{|c_a|}{\sqrt{w_a}}\le1+C\delta
 \quad(a\in\mathcal A).
 \end{gathered}
\end{equation}
The coefficient ratios in \eqref{eq:balanced} are also at most $2$.
\end{proposition}

We first construct the one-dimensional factors that will spread each
original frequency into a box.  The requirement that every coefficient
have the same magnitude makes the later comparison with $w_a$ exact.

\begin{lemma}[Spreading with equal coefficient magnitudes]
\label{lem:equal-modulus-spreading}
Let $\mathcal B\subset\R\setminus\{0\}$ be finite and invariant under
negation.  For every $\alpha>0$, there is a real
$g\in C_c^\infty((-1/2,1/2))$, with $0\le g\le1$, such that for every
sufficiently large real $T$ and each $\beta\in\mathcal B$ there is a
trigonometric polynomial
\[
 f_\beta(y)=\sum_{|u|\le\lfloor T|\beta|/2\rfloor}a_{\beta,u}\e(uy),
 \qquad
 |a_{\beta,u}|=L_\beta^{-1/2},\qquad
 L_\beta=2\lfloor T|\beta|/2\rfloor+1,
\]
satisfying
\begin{equation}\label{eq:spreading-approximation}
 \norm{f_\beta-g(\,\cdot\,)\e(T\beta(\,\cdot\,)^2/2)}_\infty\le\alpha.
\end{equation}
Here the target is extended periodically from $[-1/2,1/2]$, and the
choices may satisfy $f_{-\beta}=\overline{f_\beta}$.
\end{lemma}

\begin{proof}
Choose $0<q<1/4$, to be fixed in terms of $\alpha$, and then
choose such a $g$ with $\int g^2\ge1-q$.
For a fixed $\beta\in\mathcal B$, write
\[
 I_\beta(u)=\int_{-1/2}^{1/2}g(y)\e(T\beta y^2/2-uy)\,dy
 \qquad(u\in\Z).
\]
Lemma~\ref{lem:quadratic} and $L_\beta/(T|\beta|)\to1$ give, uniformly
in $u$,
\begin{equation}\label{eq:spreading-coefficient-bound}
 \sqrt{L_\beta}|I_\beta(u)|\le1+q
\end{equation}
for all sufficiently large $T$.
The coefficients outside the stated interval have a small sum of
absolute values.  Indeed, $\supp g$ is separated from both endpoints,
so on a neighborhood of this support,
\[
 |T\beta y-u|\ge c_{g,\beta}(T+|u|)
 \quad\text{if }|u|>T|\beta|/2.
\]
Twice integrating by parts, with no boundary terms, gives
$|I_\beta(u)|\le C_{g,\beta}(T+|u|)^{-2}$ there.  The derivative of
$T\beta y-u$ is $T\beta$, so all terms from these integrations have
this bound.  Consequently
\begin{equation}\label{eq:spreading-tail}
 \sum_{|u|>\lfloor T|\beta|/2\rfloor}|I_\beta(u)|=O_{g,\beta}(T^{-1}).
\end{equation}

For the retained indices define
$Z_u=\sqrt{L_\beta}I_\beta(u)/(1+q)$.
By \eqref{eq:spreading-coefficient-bound}, $|Z_u|\le1$.
Parseval's identity and \eqref{eq:spreading-tail} show that, for large $T$,
\[
 \frac1{L_\beta}\sum_u|Z_u|^2\ge\frac{1-2q}{(1+q)^2}.
\]
Since $1-|Z_u|\le1-|Z_u|^2$, the defect
$\mu=\frac12\sum_u(1-|Z_u|)$ satisfies
\[
 \frac{\mu}{L_\beta}\le\frac{4q+q^2}{2(1+q)^2}\le2q<1/2.
\]
Apply Lemma~\ref{lem:complex-rounding} on this consecutive interval of
frequencies, and divide its conclusion by $\sqrt{L_\beta}$.  It gives
unit complex numbers $\eta_u$.  Their polynomial
$f_\beta(y)=\sum_u L_\beta^{-1/2}\eta_u\e(uy)$ satisfies
\[
 \norm{f_\beta-\frac1{1+q}\sum_u I_\beta(u)\e(u\,\cdot)}_\infty
 \le C\left(L_\beta^{-1/2}+\sqrt{q\log(80/q)}\right).
\]
Indeed, for $x=\mu/L_\beta>0$, monotonicity gives
$x\log(80/x)\le2q\log(40/q)\le2q\log(80/q)$; the zero-defect
case follows directly from the rounding lemma.  By \eqref{eq:spreading-tail}, the polynomial before rounding
differs from the target by at most $q+o(1)$ in supremum norm.
First choose $q$ small enough, then fix $g$, and finally take $T$ large.
The resulting error is at most $\alpha$.
All estimates hold simultaneously for the finite set $\mathcal B$.
Construct the polynomials for positive $\beta$ and set
$f_{-\beta}=\overline{f_\beta}$: the targets are conjugates and the
frequency intervals are symmetric, so all conclusions are preserved.
\end{proof}

\begin{proof}[Proof of Proposition~\ref{prop:balanced}]
We first construct a real polynomial with norm at most one and squared
$L^2$ norm close to one.  Starting with $p_0=0$, define successively
\[
 p_j(y_1,\ldots,y_j)
 =p_{j-1}(y_1,\ldots,y_{j-1})
  +\tfrac12(1-p_{j-1}(y_1,\ldots,y_{j-1})^2)\cos(2\pi y_j).
\]
These polynomials have mean zero and satisfy $|p_j|\le1$, because
$(1-r^2)/2\le1-|r|$ for $|r|\le1$.
Writing $M_j=\int p_j^2$, integration in the new coordinate gives
\[
 M_j-M_{j-1}=\tfrac18\int(1-p_{j-1}^2)^2
 \ge\tfrac18(1-M_{j-1})^2.
\]
Thus $M_j\uparrow1$: a limit smaller than one would force increments
bounded away from zero.  Fix $d\ge2$ with $M_d\ge1-\delta$, and write
\[
 p_d(y)=\sum_{s\in\mathcal S}b_s\e(s\cdot y),
 \qquad b_s\ne0.
\]
Zero is not in $\mathcal S$, since $p_d$ has mean zero.  At each step,
old frequencies acquire a zero final coordinate and persist, while
new frequencies have final coordinate $1$ or $-1$.  Hence every
$s\in\mathcal S$ has last nonzero coordinate $1$ or $-1$.
Parallel members of $\mathcal S$ are therefore equal up to sign.
The first-coordinate pair persists, and the second step introduces
another pair, since $\int(1-p_1^2)>0$.
Choose representatives $s_1,\ldots,s_D$ of these sign pairs.

We next choose the sizes of the frequency boxes.  Choose
$\gamma\in\R^d$ outside the finitely many hyperplanes $s^\perp$, and
put $\rho_s=|s\cdot\gamma|>0$.  The box over $s$ should contain a
number of frequencies approximately proportional to $|b_s|^2/\rho_s$:
equal spreading then makes each squared coefficient approximately
proportional to $\rho_s$, matching the weight supplied by a vector $v$
proportional to $(\gamma,0)$.  We will choose vectors
$W_1,\ldots,W_D\in\R^d$ with nonzero projections $s\cdot W_j$ and use
these projections as spreading parameters.  The box over $s$ will then
have size asymptotic to $T^D\prod_j|s\cdot W_j|$.  We therefore require
these products to be approximately $\tau|b_s|^2/\rho_s$, using one
small positive parameter $\tau$ for all $s$.

For each $j$, choose $W_j^0\in s_j^\perp$ not orthogonal to any other
representative.  This is possible because $d\ge2$ and the representatives
are pairwise nonparallel, so only finitely many proper subspaces of
$s_j^\perp$ are excluded.  Choose $W_j^1\in\R^d$ such that
\[
 s_j\cdot W_j^1
 =\frac{|b_{s_j}|^2}
 {\rho_{s_j}\prod_{\ell\ne j}|s_j\cdot W_\ell^0|}.
\]
Set $W_j=W_j^0+\tau W_j^1$.  As $\tau\downarrow0$,
\[
 \frac1\tau\prod_{j=1}^D|s\cdot W_j|
 \longrightarrow\frac{|b_s|^2}{\rho_s}
 \qquad(s\in\mathcal S).
\]
Indeed, exactly the factor belonging to the representative of $\{s,-s\}$
vanishes at $\tau=0$, and its first-order coefficient was just prescribed.
Fix $\tau>0$ small enough that all factors are nonzero and, with
$\beta_{s,j}=s\cdot W_j$ and $\Delta_s=\prod_j|\beta_{s,j}|$,
\begin{equation}\label{eq:projection-products}
 (1-\delta)\frac{\tau|b_s|^2}{\rho_s}
 \le\Delta_s\le
 (1+\delta)\frac{\tau|b_s|^2}{\rho_s}
 \qquad(s\in\mathcal S).
\end{equation}

All these data are now fixed.  Apply
Lemma~\ref{lem:equal-modulus-spreading} to the finite list of nonzero
$\beta_{s,j}$, choosing $\alpha>0$ small enough that
\begin{equation}\label{eq:balanced-product-error}
 \left(\sum_{s\in\mathcal S}|b_s|\right)
 D\alpha(1+\alpha)^{D-1}\le\delta.
\end{equation}
With the resulting common cutoff $g$ and sufficiently large $T$, define
\[
 F(y,z)=\sum_{s\in\mathcal S}b_s\e(s\cdot y)
                  \prod_{j=1}^D f_{\beta_{s,j}}(z_j),
 \qquad (y,z)\in\T^d\times\T^D.
\]
The conjugation condition on the factors and
$b_{-s}=\overline{b_s}$ show that $F$ is real.
If each factor is replaced by its target, the resulting function is
\[
 \left(\prod_{j=1}^Dg(z_j)\right)
 p_d\left(y+\frac T2\sum_{j=1}^D W_jz_j^2\right),
\]
where $z_j$ is represented in $[-1/2,1/2]$.  The cutoffs vanish near
the endpoints, so this expression is well defined on the torus and has
norm at most one.  A telescoping product estimate bounds its distance
from $F$ by the left side of \eqref{eq:balanced-product-error}.
Therefore $\norm F_\infty\le1+\delta$.

Over each old frequency $s$, the Fourier support of $F$ is the full box
\[
 a=(s,u_1,\ldots,u_D),\qquad
 |u_j|\le\lfloor T|\beta_{s,j}|/2\rfloor.
\]
This box has $n_s=\prod_jL_{\beta_{s,j}}$ members, each with coefficient
magnitude $|b_s|/\sqrt{n_s}$.  Distinct old frequencies give disjoint
boxes, so every indicated coefficient is nonzero.  If two resulting
frequencies are parallel, their $s$ coordinates force the proportionality
factor to be $1$ or $-1$; thus they are equal or opposite.  There remain
at least two sign pairs, and no zero frequency.

Since $L_\beta/(T|\beta|)\to1$, increase $T$ so that
\begin{equation}\label{eq:box-sizes}
 (1-\delta)T^D\Delta_s\le n_s\le(1+\delta)T^D\Delta_s
 \qquad(s\in\mathcal S),
\end{equation}
and choose the approximants at this value of $T$.  This is permitted
because the spreading lemma holds for every sufficiently large $T$.
Fix $T$, put $m=d+D$, and take
\[
 v=\left(\frac{\gamma}{\tau T^D(1+\delta)^3},0\right)\in\R^{d+D}.
\]
For every frequency $a$ over $s$, its weight is
$w_a=\rho_s/(\tau T^D(1+\delta)^3)>0$.
Combining \eqref{eq:projection-products} and \eqref{eq:box-sizes} gives
\begin{align}
 \frac{(1-\delta)^2}{(1+\delta)^3}|b_s|^2
 &\le\sum_{a\text{ over }s}w_a\le\frac{|b_s|^2}{1+\delta},
 \label{eq:balanced-weight-bounds}\\
 1+\delta&\le\frac{|c_a|^2}{w_a}
 \le\frac{(1+\delta)^3}{(1-\delta)^2}.
 \label{eq:balanced-ratio-bounds}
\end{align}
Finally, $1-\delta\le\sum_s|b_s|^2=M_d\le1$, so the total weight lies
between $(1-\delta)^3/(1+\delta)^3$ and $1/(1+\delta)$.
The ratio bounds in \eqref{eq:balanced-ratio-bounds} imply
$1\le|c_a|/\sqrt{w_a}\le1+C\delta$, and the upper bound is less than
$2$ for $0<\delta<1/100$.  These are all the assertions of
\eqref{eq:balanced}.  The choices were made in the order
$\delta$, $p_d$, $\gamma$, $(W_j^0,W_j^1)$, $\tau$, $\alpha$, $g$, $T$;
in particular the entire auxiliary construction is fixed before the
polynomial length is chosen.
\end{proof}
